\documentclass[12pt]{article}
\usepackage{amsmath, amssymb}
\usepackage[margin=1in]{geometry}
\setlength{\parskip}{6pt}
\title{QUBO Formulation: Network Community Detection Selection}
\date{}
\begin{document}
\maketitle

\subsection*{Network Community Detection Selection}

\paragraph{Problem Statement.}
Given a set of $M$ candidate community partitions for a network, each associated with a pre-computed quality score $s_c$, select exactly one partition that maximizes the total quality score. The optimization problem is a discrete selection task over a finite candidate set.

\paragraph{Decision Variables.}
For each candidate partition $c \in \{0, 1, \dots, M-1\}$, introduce a binary decision variable $x_c \in \{0, 1\}$, where $x_c = 1$ indicates that partition $c$ is selected, and $x_c = 0$ otherwise. The vector of decision variables is denoted $\mathbf{x} = (x_0, x_1, \dots, x_{M-1})^\top$.

\paragraph{QUBO Derivation.}
\textbf{Step 1 --- Original Objective.} The goal is to maximize the total quality score $\sum_{c=0}^{M-1} s_c x_c$. Quadratic unconstrained binary optimization (QUBO) solvers minimize energy, so we negate the objective:
\begin{align}
\min_{\mathbf{x} \in \{0,1\}^M} \; -\sum_{c=0}^{M-1} s_c x_c.
\end{align}

\textbf{Step 2 --- Constraint Formulation.} Exactly one partition must be selected, which imposes the equality constraint:
\begin{align}
\sum_{c=0}^{M-1} x_c = 1.
\end{align}

\textbf{Step 3 --- Penalty Term Expansion.} We enforce the constraint by adding a squared penalty term with weight $\lambda > 0$:
\begin{align}
P(\mathbf{x}) = \lambda \left( \sum_{c=0}^{M-1} x_c - 1 \right)^2.
\end{align}
Expanding the square yields:
\begin{align}
P(\mathbf{x}) = \lambda \left( \sum_{c=0}^{M-1} x_c^2 + \sum_{c=0}^{M-1} \sum_{d=0, d \neq c}^{M-1} x_c x_d - 2 \sum_{c=0}^{M-1} x_c + 1 \right).
\end{align}
Applying the idempotent property $x_c^2 = x_c$ for binary variables simplifies the expression to:
\begin{align}
P(\mathbf{x}) = \lambda \left( 1 - \sum_{c=0}^{M-1} x_c + \sum_{c=0}^{M-1} \sum_{d=0, d \neq c}^{M-1} x_c x_d \right).
\end{align}
Distributing $\lambda$ gives the general penalty contribution:
\begin{align}
P(\mathbf{x}) = \lambda - \lambda \sum_{c=0}^{M-1} x_c + \lambda \sum_{c=0}^{M-1} \sum_{d=0, d \neq c}^{M-1} x_c x_d.
\end{align}

\textbf{Step 4 --- Combined Hamiltonian.} Adding the negated objective and the penalty term produces the total QUBO Hamiltonian $H(\mathbf{x})$:
\begin{align}
H(\mathbf{x}) = -\sum_{c=0}^{M-1} s_c x_c + \lambda - \lambda \sum_{c=0}^{M-1} x_c + \lambda \sum_{c=0}^{M-1} \sum_{d=0, d \neq c}^{M-1} x_c x_d.
\end{align}
Grouping linear and quadratic terms:
\begin{align}
H(\mathbf{x}) = \lambda + \sum_{c=0}^{M-1} (-s_c - \lambda) x_c + \lambda \sum_{c=0}^{M-1} \sum_{d=0, d \neq c}^{M-1} x_c x_d.
\end{align}
To match the standard symmetric QUBO form $H(\mathbf{x}) = c + \sum_{i} Q_{ii} x_i + \sum_{i < j} 2 Q_{ij} x_i x_j$, we note that $\sum_{c \neq d} x_c x_d = 2 \sum_{c < d} x_c x_d$. Substituting this identity yields:
\begin{align}
H(\mathbf{x}) = \lambda + \sum_{c=0}^{M-1} (-s_c - \lambda) x_c + \sum_{0 \leq c < d \leq M-1} 2\lambda x_c x_d.
\end{align}

\textbf{Step 5 --- Q Matrix and Offset.} Reading off the coefficients from the standard form gives the closed-form expressions for the QUBO parameters:
\begin{align}
Q_{cc} &= -s_c - \lambda, & \forall c \in \{0, \dots, M-1\}, \\
Q_{cd} &= \lambda, & \forall 0 \leq c < d \leq M-1, \\
c &= \lambda.
\end{align}

\paragraph{Penalty Weight Justification.}
The penalty weight $\lambda$ must be chosen sufficiently large to ensure that any violation of the equality constraint incurs a higher energy cost than the maximum possible gain from the objective function. The maximum possible objective value achievable by any feasible solution is $\max_{c} s_c$. If the constraint is violated, the term $\left( \sum_{c} x_c - 1 \right)^2$ evaluates to at least $1$ for integer variables, contributing at least $\lambda$ to the energy. Therefore, selecting $\lambda > \max_{c} s_c$ guarantees that the energy of any infeasible solution strictly exceeds the energy of any feasible solution.

\paragraph{Correctness Argument.}
Minimizing $H(\mathbf{x})$ over $\{0,1\}^M$ recovers the optimal partition because feasible solutions (where exactly one $x_c = 1$) incur zero penalty, so their energy equals the negative of the objective value. Infeasible solutions violate the constraint, incurring a penalty of at least $\lambda$. By the choice of $\lambda$, the energy of any infeasible solution is strictly greater than the energy of any feasible solution. Consequently, the global minimum must be feasible, and among feasible solutions, the minimum energy corresponds to the maximum objective value, i.e., the partition with the highest quality score.

\paragraph{Illustrative Example.}
Consider an instance with $M = 3$ candidate partitions and pre-computed scores $s_0 = 0.8$, $s_1 = 0.5$, and $s_2 = 0.9$. We choose $\lambda = 1.0$, which satisfies $\lambda > \max\{0.8, 0.5, 0.9\}$. Substituting these values into the general formulas yields the concrete QUBO parameters:
\begin{align}
Q_{00} &= -0.8 - 1.0 = -1.8, & Q_{11} &= -0.5 - 1.0 = -1.5, & Q_{22} &= -0.9 - 1.0 = -1.9, \\
Q_{01} &= 1.0, & Q_{02} &= 1.0, & Q_{12} &= 1.0, \\
c &= 1.0.
\end{align}
The resulting Hamiltonian is:
\begin{align}
H(x_0, x_1, x_2) = 1.0 - 1.8 x_0 - 1.5 x_1 - 1.9 x_2 + 2.0 x_0 x_1 + 2.0 x_0 x_2 + 2.0 x_1 x_2.
\end{align}
Evaluating $H$ over the three feasible assignments:
\begin{align}
H(1, 0, 0) &= 1.0 - 1.8 = -0.8, \\
H(0, 1, 0) &= 1.0 - 1.5 = -0.5, \\
H(0, 0, 1) &= 1.0 - 1.9 = -0.9.
\end{align}
The global minimum is $-0.9$ at $\mathbf{x} = (0, 0, 1)^\top$, which correctly selects candidate partition $2$ with the highest score $s_2 = 0.9$.

\end{document}